% Reproduced inherited proof; finding P-F30.
\subsection{The full cubic support theorem}

\begin{theorem}[P-F30: extremal two-power determinant character forces the chain]
\label{thm:fullcubicsupport}
Let $n\ge2$, let $F\in\mathbb C[x_1,\ldots,x_n]$ be
a cubic, and let $B$ be the matrix of exponent vectors
of its nonzero monomials. Suppose
$\mathbf1=(1,\ldots,1)^t\in\operatorname{im}_{\mathbb Q}B$.
If a diagonal matrix $A$ of finite two-power order
fixes $F$ exactly, then
\[
 \operatorname{ord}(\det A)\le2^{n-1}.
 \tag{P-F30a}
\]
If equality holds, then, after a permutation of variables,
the support of $F$ is exactly
\[
 x_1^2x_2,\ x_2^2x_3,\ldots,\ x_{n-1}^2x_n,\ x_n^3.
 \tag{P-F30b}
\]
Its coefficients can be made one by a diagonal substitution,
and $A$ itself has order $2^{n-1}$. In particular, the
form is smooth. The bound is attained by the chain.

The rational-span hypothesis holds for every
$SL_n$-semistable cubic. Thus this theorem applies
to every eigenbasis of the finite lifts in P-F26;
it imposes no prior bound on the order of the lift.
\end{theorem}

\begin{proof}
The exponent character
\[
 \psi:\mathbb Z^n\longrightarrow\mathbb C^*,
 \qquad (a_i)\longmapsto\prod_i\alpha_i^{a_i},
\]
where $\alpha_i$ are the eigenvalues of $A$, annihilates
the support lattice and has two-power image. Its value
on $\mathbf1$ is $\det A$.
We will use the following elementary denominator fact
throughout. If $K$ is a selection of support columns
forming a basis of the rational support span, then
$Kz=\mathbf1$ has a unique rational solution $z$.
For any nonzero maximal square minor $K_I$, Cramer's
rule shows that the least positive integer $d$ for
which $dz$ is integral divides $|\det K_I|$.
Since $d\mathbf1=K(dz)$ is in the support lattice,
the order of $\det A$ divides $d$. Consequently
\[
 \operatorname{ord}(\det A)\le2^{v_2(\det K_I)}.
 \tag{P-F30c}
\]
This argument requires neither full rank of $B$
nor a bound on the orders of the individual $\alpha_i$.

Form the graph whose edges are the supported monomials
$x_i^2x_j$, with coefficient two at $i$ and one at $j$,
and whose loops are the supported cubes. Include all
$n$ vertices, even those isolated in this graph.
Let $G$ be the matrix of these graph-support columns.
Every other support column has the form
$e_i+e_j+e_k$ for three distinct vertices.

In a connected graph component on $v$ vertices,
a spanning-tree matrix has rank $v-1$. Thus the
graph-component matrix has rank $v$ or $v-1$.
This includes an isolated vertex: it has rank one
if its cube occurs, and rank zero otherwise.
Write $f$ for the number of full-rank components,
$c$ for the number of rank-deficient components,
and $v_i$ for their vertex counts. Then
\[
 \operatorname{rank}G=n-c.
 \tag{P-F30d}
\]

Choose a square support basis $C_i$ in each full-rank
component by adjoining one independent graph column
to a spanning tree. The unicyclic determinant
calculation in P-F27 gives
\[
 v_2(\det C_i)\le v_i-1.
 \tag{P-F30e}
\]
For precision, leaf expansion contributes a factor
one or two for each edge off the unique cycle.
A loop contributes the odd factor three.
A cycle of length $l\ge2$ contributes, up to sign,
$2^a+(-1)^{l-1}2^{l-a}$; if nonzero, its valuation
is at most $\lfloor(l-1)/2\rfloor$. These alternatives
give (P-F30e), including loops and double edges.

Choose a spanning-tree matrix $T_i$ in each deficient
component. Its signed cofactor vector has the form
\[
 2^{e_i}u_i,\qquad
 (u_i)_v=\varepsilon_i(-2)^{h_{iv}},
 \quad \varepsilon_i\in\{1,-1\},\quad
 \min_v h_{iv}=0.
 \tag{P-F30f}
\]
All exponents here are nonnegative integers.
Indeed every maximal tree minor is, up to sign,
a power of two, by successive leaf expansion.
The cofactor relation along an edge with
coefficients two and one is $2u_v+u_w=0$.
Thus adjacent heights differ by one and their
signs are opposite; connectivity gives the
single sign convention in (P-F30f).
For a single isolated deficient vertex, the
empty minor is one and we use $e_i=h_{iv}=0$.

Put
\[
 b_i=v_i-1-e_i,\qquad H_i=\max_v h_{iv}.
\]
Since a tree minor is the product of $v_i-1$
factors, each one or two, one has
\[
 0\le H_i\le b_i.
 \tag{P-F30g}
\]
We need one additional equality property: if
$H_i=b_i$, the vertex of height $H_i$ is unique.
To prove it, choose such a vertex as root.
Its deleted-row minor has absolute value
$2^{e_i+H_i}=2^{v_i-1}$. Expand this minor by
removing leaves toward that root. Every factor
must be two. Thus each edge has coefficient
two at its child and one at its parent.
The cofactor relation gives
$h_{\rm parent}=h_{\rm child}+1$.
Heights strictly decrease away from the root,
so no other vertex has its height.
The assertion is also immediate for an isolated vertex.

Normalize the entries in each of these cofactor
vectors by $2^{b_i}$. For sums over distinct
vertices of one deficient component, the following
bounds hold:
\[
 \left|\sum_{v\in E}\frac{(u_i)_v}{2^{b_i}}\right|
 \le
 \begin{cases}
  1,&|E|=1,\\
  5/4,&|E|=2,\\
  3/2,&|E|=3.
 \end{cases}
 \tag{P-F30h}
\]
Here a case is used only when the component has
enough distinct vertices. If $H_i<b_i$, every
entry has absolute value at most $1/2$, which
proves all three bounds. If $H_i=b_i$, there is
only one entry of absolute value one. Change
the overall sign to make this entry positive.
Every entry of absolute value $1/2$ is negative,
because its height is $H_i-1$; every remaining
entry has absolute value at most $1/4$.
A pair avoiding the top entry has absolute sum
at most one; a pair containing it has absolute
sum at most $5/4$. A triple avoiding it has
absolute sum at most $3/2$. For a triple containing
it, either the other two entries are both $-1/2$,
giving zero, exactly one is $-1/2$, giving absolute
sum at most $3/4$, or neither is, giving at most
$1+1/4+1/4=3/2$. This proves (P-F30h).

The chosen full-component bases and deficient-component
trees together form a rational basis for the columns
of $G$. Extend it to a basis for all support columns
by adjoining $q$ independent triple columns
$\beta_1,\ldots,\beta_q$. Then
\[
 0\le q\le c,\qquad
 r:=\operatorname{rank}B=n-c+q.
 \tag{P-F30i}
\]
If $q=0$, the graph subform itself satisfies the
rational-span hypothesis. P-F27 applies directly:
it gives (P-F30a), and equality gives the chain
eigenvalue list and the chain graph support.
That eigenvalue list admits no other cubic
monomials, so there can be no additional triple
monomials in $F$. This proves all the required
claims in the case $q=0$.
We therefore assume $q>0$ for the rest of the argument.

Define the $c\times q$ compressed matrix
\[
 D_{ij}=u_i^t(\beta_j|_i),
 \tag{P-F30j}
\]
where $\beta_j|_i$ is the restriction to the vertices
of deficient component $i$. The map on the entire
vertex space that sends a vector to these $c$
cofactor pairings is surjective, and its kernel
is precisely the rational graph-column span:
on a full component the graph columns span all
coordinates, while on a deficient component the
kernel of $u_i^t$ is the tree-column span.
Consequently $D$ has rank $q$. Choose $q$
deficient-component indices $J$ for which the
square matrix $D_J$ has nonzero determinant.

Let $K$ be the $n\times r$ original support-basis
matrix just selected. We now specify an actual
nonzero $r\times r$ minor of it. Keep every row
in each full component and in each deficient
component indexed by $J$. In every other deficient
component, delete one row at a vertex of height zero,
and keep the remaining rows. The resulting row
set $I$ has size $n-(c-q)=r$. Its determinant is
\[
 \det K_I
   =\pm\left(\prod_{\text{full }i}\det C_i\right)
          2^{\sum_{\text{deficient }i}e_i}
          \det D_J.
 \tag{P-F30k}
\]
We justify this factorization explicitly.
In a full component the graph basis is square;
eliminating the triple-column entries there
using its graph columns and expanding factors
out $\det C_i$. In a deficient component not
in $J$, the retained tree block is square.
Its deleted row had height zero, so its determinant
is $\pm2^{e_i}$; eliminate the triple entries
there and factor out that determinant as well.
These column eliminations do not alter other
components, since every graph column is supported
in its own component.
The remaining matrix consists of the tall
tree blocks for $i\in J$ and $q$ triple columns.
Expansion by the maximal minors of each tree
block leaves one row pairing for each such
component. Its signed cofactor vector is
$2^{e_i}u_i$. The remaining $q\times q$
determinant is exactly the pairings (P-F30j),
giving the remaining factors
$2^{\sum_{i\in J}e_i}\det D_J$.
This proves (P-F30k), and proves that the
specified original minor is nonzero.

Divide row $i$ of $D_J$ by $2^{b_i}$.
Each resulting column is formed from at most
three distinct vertices of the original graph.
If they belong to three selected components,
its squared Euclidean norm is at most three.
If two lie in one selected component and one
in another, (P-F30h) gives squared norm at most
\[
 (5/4)^2+1=41/16<3.
\]
If all three lie in one component, the squared
norm is at most $(3/2)^2<3$. Ignoring vertices
in unselected or full components only leaves
cases with fewer than three contributing vertices,
which satisfy the same bound. Every normalized
column therefore has Euclidean norm at most
$\sqrt3$. Hadamard's determinant inequality yields
\[
 0<|\det D_J|
       \le 2^{\sum_{i\in J}b_i}3^{q/2}.
 \tag{P-F30l}
\]

Suppose first that $q<c$, or that $q=c$ and $f>0$.
Using (P-F30e), (P-F30k), and (P-F30l) gives
\begin{align*}
 v_2(\det K_I)
 &\le \sum_{\text{full }i}(v_i-1)
      +\sum_{\text{deficient }i}e_i
      +\sum_{i\in J}b_i+\frac q2\log_2 3\\
 &=n-f-c-\sum_{\substack{\text{deficient }i\\i\notin J}}b_i
              +\frac q2\log_2 3\\
 &<n-f-c+q.
 \tag{P-F30m}
\end{align*}
The last inequality is strict because $q>0$
and $\sqrt3<2$; the omitted $b_i$ are nonnegative.
If $q<c$, its right side is at most $n-1$.
If $q=c$ and $f>0$, it is likewise at most $n-1$.
The valuation is an integer, so in both cases
\[
 v_2(\det K_I)\le n-2.
 \tag{P-F30n}
\]

The only case remaining is $q=c$ and $f=0$.
Here $K$ is a square $n\times n$ support matrix,
and $I$ is its full row set. Every original
support column has sum three. Hence
$\mathbf1^tK=3\mathbf1^t$, and reduction modulo
three shows that $3$ divides its nonzero determinant.
By (P-F30k), which now has no full-component
factors, $3$ also divides $\det D$.
The two-part of this determinant is consequently
at most $|\det D|/3$. Since
$\sum_i(e_i+b_i)=\sum_i(v_i-1)=n-c$,
we obtain
\begin{align*}
 2^{v_2(\det K)}
 &=2^{\sum_i e_i}\,2^{v_2(\det D)}\\
 &\le 2^{\sum_i e_i}\frac{|\det D|}{3}\\
 &\le 2^{n-c}3^{c/2-1}
  <\frac{2^n}{3}
  <2^{n-1}.
 \tag{P-F30o}
\end{align*}
Here $c=q>0$ and $3^{c/2}<2^c$ justify the
first strict inequality. Thus (P-F30n) holds
in this case as well.

Combining (P-F30c) and (P-F30n) proves the
stronger bound
$\operatorname{ord}(\det A)\le2^{n-2}$
whenever $q>0$. All extremal cases therefore
occur in the already treated case $q=0$,
where P-F27 proves precisely (P-F30b), the
coefficient normalization, and the exact order
of $A$. Conversely the chain is fixed by
the diagonal matrix with eigenvalues
$\xi,\xi^{-2},\ldots,\xi^{(-2)^{n-1}}$
for a primitive $2^{n-1}$th root $\xi$.
Its determinant exponent
$1-2+4-\cdots+(-2)^{n-1}$ is odd, so its
determinant has order $2^{n-1}$. The bound
is therefore attained.

Finally, P-F27's coefficient-monomial argument
shows that any positive-degree $SL_n$ invariant
nonzero at $F$ expresses $\mathbf1$ as a
nonnegative rational combination of the
support exponents. This supplies the
rational-span hypothesis for every semistable
cubic, in every coordinate basis, and completes
the proof.
\end{proof}

\paragraph{The structural conclusion.}
The theorem closes the mixed-support case left by
P-F27 and P-F29. It uses the actual determinant
character, rather than identifying that character's
order with an arbitrary lift order. Its proof
allows rank-deficient exponent matrices and uses
an explicitly constructed maximal minor in those
cases. Combining it with P-F26 rules out
nonzero singular representatives in the
$\boldsymbol\mu_{2^{n-1}}$ fixed quotient locus.
The ensuing global fixed-locus and Hilbert-pole
conclusions are stated separately.
